%%%%%%%%%%%%%%%%%%%%%%%%%%%%%%%%%%%%%%%%%%%%%%%%%%%%%%%%%%%%%%%%%%%%%%%%%%%%%%%%%%
% NBA case study, extracted from ml_knn_sklearn.tex (Sep 2026): the only genuinely
% unique content in that file for MLCoEP purposes -- the Iris/sklearn walkthrough
% earlier in that file stays excluded here since it overlaps Session 7's sklearn
% coverage (ml_intro_sklearn.tex), but this hands-on distance/regression exercise
% is not covered anywhere else in the session.
\begin{frame}[fragile]\frametitle{}
\begin{center}
{\Large Test case: NBA Players}
\end{center}
\end{frame}

%%%%%%%%%%%%%%%%%%%%%%%%%%%%%%%%%%%%%%%%%%%%%%%%%%%%%%%%%%%
\begin{frame}[fragile]\frametitle{Exercise: Loading the NBA Dataset}
A look at the data. Each row contains information on how a player performed in the
2013-2014 NBA season.
\begin{lstlisting}
import pandas
with open("nba_2013.csv", 'r') as csvfile:
    nba = pandas.read_csv(csvfile)

# The names of all the columns in the data.
print(nba.columns.values)
\end{lstlisting}
\end{frame}

%%%%%%%%%%%%%%%%%%%%%%%%%%%%%%%%%%%%%%%%%%%%%%%%%%%%%%%%%%%
\begin{frame}[fragile]\frametitle{Exercise: Dataset Columns}
Here are some selected columns from the data:
\begin{itemize}
\item     player - name of the player
\item     pos - the position of the player
\item     g - number of games the player was in
\item     gs - number of games the player started
\item     pts - total points the player scored
\end{itemize}
There are many more columns in the data, mostly containing information about average player game performance over the course of the season.
\end{frame}

%%%%%%%%%%%%%%%%%%%%%%%%%%%%%%%%%%%%%%%%%%%%%%%%%%%%%%%%%%%
\begin{frame}[fragile]\frametitle{Exercise: Distance to LeBron James}
Study distance: find the most similar NBA players to Lebron James.
\begin{lstlisting}
import math

# Extract LeBron's row to use as the reference point
selected_player = nba[nba["player"] == "LeBron James"].iloc[0]

# Numeric columns to measure player similarity across stats
distance_columns = ['age', 'g', 'gs', 'mp', 'fg', 'fga',
						'fg.', 'x3p', 'x3pa', 'x3p.', 'x2p',
						 'x2pa', 'x2p.', 'efg.', 'ft', 'fta',
						  'ft.', 'orb', 'drb', 'trb', 'ast', 'stl',
						  'blk', 'tov', 'pf', 'pts']

def euclidean_distance(row):
    # Sum squared differences across all stat columns vs LeBron
    inner_value = 0
    for k in distance_columns:
        inner_value += (row[k] - selected_player[k]) ** 2
    return math.sqrt(inner_value)

# Compute distance from every player to LeBron row-by-row
lebron_distance = nba.apply(euclidean_distance, axis=1)
\end{lstlisting}
\end{frame}

%%%%%%%%%%%%%%%%%%%%%%%%%%%%%%%%%%%%%%%%%%%%%%%%%%%%%%%%%%%
\begin{frame}[fragile]\frametitle{Exercise: Normalizing Columns}
Normalizing columns
\begin{lstlisting}
# Select only the numeric columns from the NBA dataset
nba_numeric = nba[distance_columns]

# Normalize all of the numeric columns
nba_normalized = (nba_numeric - nba_numeric.mean()) / nba_numeric.std()


# Fill in NA values in nba_normalized
nba_normalized.fillna(0, inplace=True)
\end{lstlisting}
\end{frame}

%%%%%%%%%%%%%%%%%%%%%%%%%%%%%%%%%%%%%%%%%%%%%%%%%%%%%%%%%%%
\begin{frame}[fragile]\frametitle{Exercise: Distance via SciPy}
Use the distance.euclidean function from scipy.spatial, a much faster way to calculate euclidean distance.
\begin{lstlisting}
from scipy.spatial import distance

# Find the normalized vector for lebron james.
lebron_normalized = nba_normalized[nba["player"] == "LeBron James"]

# Find the distance between lebron james and everyone else.
euclidean_distances = nba_normalized.apply(lambda row: distance.euclidean(row, lebron_normalized), axis=1)

# Create a new dataframe with distances.
distance_frame = pandas.DataFrame(data={"dist": euclidean_distances, "idx": euclidean_distances.index})
distance_frame.sort_values("dist", inplace=True)
# Find the most similar player to lebron (the lowest distance to lebron is lebron, the second smallest is the most similar non-lebron player)
second_smallest = distance_frame.iloc[1]["idx"]
most_similar_to_lebron = nba.loc[int(second_smallest)]["player"]
\end{lstlisting}
\end{frame}

%%%%%%%%%%%%%%%%%%%%%%%%%%%%%%%%%%%%%%%%%%%%%%%%%%%%%%%%%%%
\begin{frame}[fragile]\frametitle{Exercise: Train/Test Split}
Generating training and testing sets
\begin{lstlisting}
import random
from numpy.random import permutation

# Randomly shuffle the index of nba.
random_indices = permutation(nba.index)

# Set a cutoff for how many items we want in the test set (in this case 1/3 of the items)
test_cutoff = math.floor(len(nba)/3)

# Generate the test set by taking the first 1/3 of the randomly shuffled indices.
test = nba.loc[random_indices[1:test_cutoff]]

# Generate the train set with the rest of the data.
train = nba.loc[random_indices[test_cutoff:]]
\end{lstlisting}
\end{frame}

%%%%%%%%%%%%%%%%%%%%%%%%%%%%%%%%%%%%%%%%%%%%%%%%%%%%%%%%%%%
\begin{frame}[fragile]\frametitle{Exercise: KNN Regression with Scikit-Learn}
Using sklearn for k nearest neighbors
\begin{lstlisting}
x_columns = ['age', 'g', 'gs', 'mp', 'fg', 'fga', 'fg.', 'x3p', 'x3pa', 'x3p.', 'x2p', 'x2pa', 'x2p.', 'efg.', 'ft', 'fta', 'ft.', 'orb', 'drb', 'trb', 'ast', 'stl', 'blk', 'tov', 'pf']

y_column = ["pts"]

from sklearn.neighbors import KNeighborsRegressor

knn = KNeighborsRegressor(n_neighbors=5)

knn.fit(train[x_columns], train[y_column])

predictions = knn.predict(test[x_columns])
\end{lstlisting}
\end{frame}

%%%%%%%%%%%%%%%%%%%%%%%%%%%%%%%%%%%%%%%%%%%%%%%%%%%%%%%%%%%
\begin{frame}[fragile]\frametitle{Exercise: Computing Error (MSE)}
Computing error
\begin{lstlisting}
actual = test[y_column]

mse = (((predictions - actual) ** 2).sum()) / len(predictions)
\end{lstlisting}
\end{frame}

%%%%%%%%%%%%%%%%%%%%%%%%%%%%%%%%%%%%%%%%%%%%%%%%%%%%%%%%%%%
\begin{frame}[fragile]\frametitle{NBA Case Study: Recap}
\begin{block}{Intuition}
This is the same recipe taught earlier, run on real, messy data. Normalize every stat so
no single column (like total points) dominates the distance. Compute Euclidean distance
to every other player. Then either read off the single nearest neighbor (most similar to
LeBron) or average the $k$ nearest neighbors' points scored (KNN regression) to predict a
new player's performance.
\end{block}
\end{frame}
